\centerline{\bf \Huge Алгоритм Евклида}

Докажем, что $\text{НОД}(a,b)=\text{НОД}(a-b,b)$.

Покажем, что все общие делители чисел $a$ и $b$ и чисел $a-b$ и $b$ совпадают: пусть $d|a$ и $d|b$, тогда $d|(a-b)$. Также, если $d|(a-b)$ и $d|b$, тогда $d|((a-b)+b)$.

Из этих утверждений получается, что множества общих делителей чисел $a,b$ и чисел $(a-b),b$ совпадают, значит, совпадают и их наибольшие элементы. 

Повторив вычитание много раз, дойдем до остатка от деления $a$ на $b$. Получается, что  $\text{НОД}(a,b)=\text{НОД}(r,b)$, где $r$ - остаток от деления числа $a$ на $b$.

Алгоритм Евклида заключается в том, чтобы повторить операцию описанную выше много раз: 

$(a,b)=(r_1,b)=(r_1,r_2)=\ldots=(r_{n-1},r_n)$.

Где

$a=b\cdot q_1+r_1$

$b=r_1 \cdot q_2 + r_2$

$r_1=r_2 \cdot q_3 + r_3$

$\ldots$

$r_{n-1}=r_n \cdot q_{n+1} $

Процесс сходится, так как $|r_i| > |r_{i+1}|$, и число $|b|$ ограничено, бесконечный спуск невозможен.

\centerline{\bf \Huge Линейная представимость НОД}

Для любых целых $a$,$b$ существуют такие числа $x$ и $y$, что $ax+by=\text{НОД}(x,y)$.

Распишем алгоритм Евклида:

$(a,b)=(r_1,b)=(r_1,r_2)=\ldots=(r_{n-1},r_n)$, где 

$a=b\cdot q_1+r_1$

$b=r_1 \cdot q_2 + r_2$

$r_1=r_2 \cdot q_3 + r_3$

$\ldots$

$r_{n-1}=r_n \cdot q_{n+1} $


Выразим из первого уравнения $r_1=a-b\cdot q_1$ и подставим во второе:

$b=(a-b \cdot q_1) q_2 + r_2$. Отсюда выразим $r_2=b-(a-b\cdot q_1) \cdot q_2$.

Подставим все в третье уравнение:

$(a-b\cdot q_1)=(b-(a-b\cdot q_1) \cdot q_2) \cdot q_3 + r_3$ откуда мы можем выразить $r_3$ и так далее, вплоть до $r_n$. 


{\bf Пример}

Найдем НОД (28, 77).

$77=28\cdot2 +21$

$28=21 \cdot 1 +7$

$21=7 \cdot 3$

Выразим из первого $21=77-28\cdot2$

$28=(77-28\cdot2)\cdot 1+7$ откуда $28-(77-28\cdot 2)=7$

Отсюда получаем, что $7=28\cdot 3 - 77$.

Это и есть линейная представимость 7 через 28 и 77. 

$7=28\cdot 3 - 77$

\centerline{\bf \Huge Линейные Диофантовы уравнения} 

Уравнения в целых числах $ax+by = c$, где числа $a,b,c$ заданы, называются линейными диофантовыми уравнениями. 

Для их решения важна теорема:

Линейное диофантово уравнение $ax+by = c$, где числа $a,b,c$ заданы, имеет решение тогда и только тогда, когда $\text{НОД}(a,b) |c$.

{\bf Доказательство}

Пусть уравнение имеет решение, тогда получается, что левая часть уравнения делится нацело на $\text{НОД}(a,b)$, значит, правая часть делится нацело на $\text{НОД}(a,b)$, то есть $\text{НОД}(a,b) | c$

В обратную сторону: пусть $\text{НОД}(a,b) | c$, докажем, что уравнение имеет решение. По теореме о линейной представимости НОДа мы знаем, что для любых чисел $a$ и $b$ существуют целые числа $x_0$ и $y_0$ такие, что 

$$ax_0+by_0 = \text{НОД}(a,b)$$


Если $\text{НОД}(a,b)|c$, то $c = k \cdot \text{НОД}(a,b)$. Давайте умножим в выражении выше обе части уравнения на $k$, тогда получим:

$$ax_0 \cdot k + by_0\cdot k = k \cdot \text{НОД}(a,b)$$

$$a(x_0k)+b(y_0k)=c$$

Значит, $x=x_0k$ и $y=y_0k$ являются решениями диофантова уравнения, то есть уравнение имеет решение, что и требовалось доказать.


Отметим, что из доказательства теоремы следует алгоритм поиска частного решения линейных диофантовых уравнений:

Решите уравнение:

$28x+77y=35$.

Выше мы нашли НОД чисел 28 и 77 и даже его линейное представление:

$$7= 28 \cdot 3 - 77$$

Умножим обе части этого выражения на 5:

$$35 = 28 \cdot 15 - 77 \cdot 5$$

Откуда получаем частное решение $x=15$, $y=-5$. Давайте найдем все решения уравнения, зная частное:

Пусть $x_0,y_0$ - другое решение этого уравнения, подставим:

$$28x_0+77y_0=28 \cdot 15 - 77 \cdot 5$$

$$77\cdot (y_0+5)=28 \cdot (15-x_0)$$

$$11 \cdot (y_0+5)=4 \cdot (15-x_0)$$

Замечаем, что левая часть делится на 11, значит, правая часть должна делиться на 11, значит, $15 - x_0 = 11t$, подставим это в наше уравнение:

$y_0+5=4 \cdot t$, откуда получаем, что $y_0 = -5 + 4t$, $x_0=15 - 11t$, где $t$ - любое целое число. 

В общем случае уравнение решается следующим образом:

$$ax+by = c$$

Находим $\text{НОД}(a,b)$ при помощи алгоритма Евклида, далее линейно выражаем НОД. Если $c$ не делится на $\text{НОД}(a,b)$, то у уравнения решений нет, если же делится, то пусть $c= k \cdot \text{НОД}(a,b)$.

Берем линейное представление НОД: $ ax_0 + by_0 = \text{НОД}(a,b) $ и умножаем обе части на число $k$, 

$c = a(x_0k)+b(y_0k)$, откуда мы находим частное решение.

Общее решение выражается следующим образом:

$x = x_0k + \cfrac{b}{\text{НОД}(a,b)} \cdot t$

$y = y_0k - \cfrac{a}{\text{НОД}(a,b)} \cdot t$, где $t \in \mathbb{Z}$ , то есть любое целое число.